\subsection{Complex Numbers and Absolute Value}

\begin{definition} \label{def:complex-number-field-operations} \hlt{Complex Numbers}\\
Let $\C\defeq\R^2$. For $x=(a,b)$ and $y=(c,d)$ in $\C$:
\begin{enumerate}[label=(\roman*)]
\setlength{\itemsep}{0pt}
\item $x=y\Longleftrightarrow(a=c\land b=d)$.
\item $x+y\defeq(a+c,b+d)$.
\item $xy\defeq(ac-bd,ad+bc)$.
\end{enumerate}
Elements of $\C$ are \hlt{complex numbers}, and $i\defeq(0,1)\in\C$ is the \hlt{complex unit}.
\end{definition}

\begin{theorem} \label{thm:complex-numbers-form-field} \hlt{The Complex Numbers Form a Field}\\
Let $\C=\R^2$ with $(a,b)+(c,d)=(a+c,b+d)$ and $(a,b)(c,d)=(ac-bd,ad+bc)$.
\begin{enumerate}[label=(\roman*)]
\setlength{\itemsep}{0pt}
\item The additive and multiplicative identities are $(0,0)$ and $(1,0)$.
\item Every $x=(a,b)\neq(0,0)$ has inverse $x^{-1}=\left(\frac{a}{a^2+b^2},-\frac{b}{a^2+b^2}\right)$.
\item With these operations, $\C$ is a field.
\end{enumerate}
\begin{proof}
\begin{enumerate}[label=(\roman*)]
\setlength{\itemsep}{0pt}
\item For $x=(a,b)$, $x+(0,0)=(0,0)+x=x$ and $x(1,0)=(1,0)x=x$ by direct calculation.
\item At least one of $a,b$ is nonzero, so $a^2+b^2>0$ by Proposition \ref{prop:ordered-field-inequality-rules-analysis}(iv).\\
Set $q\defeq\left(\frac{a}{a^2+b^2},-\frac{b}{a^2+b^2}\right)$. Direct multiplication verifies $xq=qx=\left(\frac{a^2+b^2}{a^2+b^2},0\right)=(1,0)$.
\item Write $x=(a,b)$, $y=(c,d)$, and $u=(e,f)$. The coordinate formulas lie in $\R^2$, so both operations are closed.\\
Additive associativity is $(x+y)+u=(a+c+e,b+d+f)=x+(y+u)$.\\
Additive commutativity is $x+y=(c+a,d+b)=y+x$.\\
Together with (i) and $-(a,b)=(-a,-b)$, these identities make $(\C,+)$ an abelian group.\\
Multiplicative commutativity is $xy=(ac-bd,ad+bc)=(ca-db,cb+da)=yx$.\\
Multiplicative associativity is $(xy)u=(ace-adf-bcf-bde,acf+ade+bce-bdf)=x(yu)$.\\
The identity $x(y+u)=(a(c+e)-b(d+f),a(d+f)+b(c+e))=xy+xu$ proves left distributivity.\\
Right distributivity holds by commutativity. Parts (i)--(ii) and $(0,0)\neq(1,0)$ complete the field axioms.
\end{enumerate}
\hlt{Structural Commentary}: Coordinate identities verify the ring laws; $a^2+b^2>0$ constructs each inverse.\\
\hlt{Key Insight}: The relation $a^2+b^2>0$ makes conjugate-over-norm division possible.
\end{proof}
\end{theorem}

\begin{proposition} \label{prop:real-field-embedded-in-complex-field} \hlt{Real Numbers Inside Complex Numbers}\\
Let $\iota:\R\to\C$ be $\iota(a)\defeq(a,0)$.
\begin{enumerate}[label=(\roman*)]
\setlength{\itemsep}{0pt}
\item $(\forall a,b\in\R)\,[\iota(a+b)=\iota(a)+\iota(b)\land\iota(ab)=\iota(a)\iota(b)]$.
\item $\iota$ is injective, $\iota(0)=(0,0)$, and $\iota(1)=(1,0)$.
\end{enumerate}
\begin{proof}
\begin{enumerate}[label=(\roman*)]
\setlength{\itemsep}{0pt}
\item $\iota(a)+\iota(b)=(a,0)+(b,0)=(a+b,0)=\iota(a+b)$, and $\iota(a)\iota(b)=(a,0)(b,0)=(ab,0)=\iota(ab)$.
\item If $\iota(a)=\iota(b)$, then $(a,0)=(b,0)$, so $a=b$. The identities $\iota(0)=(0,0)$ and $\iota(1)=(1,0)$ hold directly.
\end{enumerate}
\hlt{Structural Commentary}: Zero second coordinates are stable under both operations; pair equality proves injectivity.\\
\hlt{Key Insight}: The map $a\mapsto(a,0)$ realizes real arithmetic on the first coordinate.
\end{proof}
\end{proposition}

\begin{convention} \label{conv:real-subfield-identification-complex-numbers} \hlt{Real Scalars Inside $\C$}\\
Henceforth identify $a\in\R$ with its image $\iota(a)=(a,0)\in\C$.
\end{convention}

\begin{proposition} \label{prop:complex-unit-and-coordinate-form} \hlt{Coordinate Form and $i^2$}\\
For $i=(0,1)\in\C$:
\begin{enumerate}[label=(\roman*)]
\setlength{\itemsep}{0pt}
\item $i^2=-1$.
\item $(\forall a,b\in\R)\,[(a,b)=a+bi]$.
\end{enumerate}
\begin{proof}
\begin{enumerate}[label=(\roman*)]
\setlength{\itemsep}{0pt}
\item $i^2=(0,1)(0,1)=(-1,0)=-1$.
\item $a+bi=(a,0)+(b,0)(0,1)=(a,0)+(0,b)=(a,b)$.
\end{enumerate}
\hlt{Structural Commentary}: The products $(0,1)^2=(-1,0)$ and $(b,0)(0,1)=(0,b)$ convert pair coordinates into $a+bi$.\\
\hlt{Key Insight}: Multiplication by $i=(0,1)$ transfers a real scalar to the second coordinate.
\end{proof}
\end{proposition}

\begin{definition} \label{def:complex-conjugate-real-imaginary-parts} \hlt{Conjugate, Real Part, and Imaginary Part}\\
Let $z=a+bi$ with $a,b\in\R$.
\begin{enumerate}[label=(\roman*)]
\setlength{\itemsep}{0pt}
\item The \hlt{complex conjugate} of $z$ is $\overline{z}\defeq a-bi$.
\item The \hlt{real part} and \hlt{imaginary part} of $z$ are $\operatorname{Re}(z)\defeq a$ and $\operatorname{Im}(z)\defeq b$.
\end{enumerate}
\end{definition}

\begin{proposition} \label{prop:complex-conjugation-identities} \hlt{Basic Conjugation Identities}\\
If $z,w\in\C$, then
\begin{enumerate}[label=(\roman*)]
\setlength{\itemsep}{0pt}
\item $\overline{z+w}=\overline{z}+\overline{w}$;
\item $\overline{zw}=\overline{z}\,\overline{w}$;
\item $z+\overline{z}=2\operatorname{Re}(z)$ and $z-\overline{z}=2i\operatorname{Im}(z)$;
\item $z\overline{z}\in[0,\infty)$; $z\neq0\Rightarrow z\overline{z}>0$, and $z=0\Rightarrow z\overline{z}=0$.
\end{enumerate}
\begin{proof}
Write $z=a+bi$ and $w=c+di$.
\begin{enumerate}[label=(\roman*)]
\setlength{\itemsep}{0pt}
\item $\overline{z+w}=(a+c)-(b+d)i=\overline{z}+\overline{w}$.
\item $\overline{zw}=(ac-bd)-(ad+bc)i=(a-bi)(c-di)=\overline{z}\,\overline{w}$.
\item $z+\overline{z}=2a=2\operatorname{Re}(z)$ and $z-\overline{z}=2bi=2i\operatorname{Im}(z)$.
\item $z\overline{z}=a^2+b^2\in[0,\infty)$. For $z\neq0$, $a^2+b^2>0$ by Proposition \ref{prop:ordered-field-inequality-rules-analysis}(iv), as at least one coordinate nonzero. For $z=0$, $z\overline{z}=0\overline{0}=0$.
\end{enumerate}
\hlt{Structural Commentary}: Coordinate sign reversal preserves both operations and turns $z\overline{z}$ into a sum of squares.\\
\hlt{Key Insight}: $z\overline{z}=a^2+b^2\geq0$.
\end{proof}
\end{proposition}

\begin{definition} \label{def:complex-absolute-value-analysis} \hlt{Absolute Value of a Complex Number}\\
For $z=a+bi\in\C$, define $\abs{z}\defeq\sqrt{z\overline{z}}=\sqrt{a^2+b^2}$ as the unique $r\in[0,\infty)$ with $r^2=a^2+b^2$.\\
For $z\neq0$, existence and uniqueness hold by Proposition \ref{prop:complex-conjugation-identities}(iv) and Theorem \ref{thm:nth-roots-positive-reals} with $n=2$.\\
For $z=0$, the unique nonnegative solution is $r=0$ by Proposition \ref{prop:ordered-field-inequality-rules-analysis}(iv).
\end{definition}

\begin{theorem} \label{thm:complex-absolute-value-identities} \hlt{Absolute Value Identities}\\
For $z,w\in\C$,
\begin{enumerate}[label=(\roman*)]
\setlength{\itemsep}{0pt}
\item $\abs{z}\geq0$; $z=0\Rightarrow\abs{z}=0$, and $z\neq0\Rightarrow\abs{z}>0$;
\item $\abs{\overline{z}}=\abs{z}$;
\item $\abs{zw}=\abs{z}\abs{w}$;
\item $\abs{\operatorname{Re}(z)}\leq\abs{z}$;
\item $\abs{z+w}\leq\abs{z}+\abs{w}$.
\end{enumerate}
\begin{proof}
For $r,s\in[0,\infty)$, $r^2\leq s^2\Rightarrow r\leq s$, since $s<r$ would imply $0<(r-s)(r+s)=r^2-s^2$.\\
For $t\in\R$ with $t\geq0$, $t=\abs{t}$ by equality of their nonnegative squares. For $t<0$, $t<0\leq\abs{t}$; thus $(\forall t\in\R)\,t\leq\abs{t}$.
\begin{enumerate}[label=(\roman*)]
\setlength{\itemsep}{0pt}
\item The relations $\abs{z}\geq0$ and $\abs{0}=0$ hold by Definition \ref{def:complex-absolute-value-analysis}. For $z\neq0$, $\abs{z}^2=z\overline{z}>0$ by Proposition \ref{prop:complex-conjugation-identities}(iv).\\
The inequality $\abs{z}>0$ holds by $\abs{z}^2>0$ and nonnegativity.
\item Since $\overline{\overline{z}}=z$ by definition, $\abs{\overline{z}}^2=\overline{z}z=z\overline{z}=\abs{z}^2$. Then $\abs{\overline{z}}=\abs{z}$ holds by preliminary square comparison in both directions.
\item $\abs{zw}^2=(zw)\overline{zw}=(z\overline{z})(w\overline{w})$ by Proposition \ref{prop:complex-conjugation-identities}(ii) and commutativity. $\abs{zw}^2=\abs{z}^2\abs{w}^2=(\abs{z}\abs{w})^2$.\\
Then $\abs{zw}=\abs{z}\abs{w}$ holds by the preliminary square comparison and nonnegativity.
\item For $z=a+bi$, $\abs{\operatorname{Re}(z)}^2=a^2\leq a^2+b^2=\abs{z}^2$. Then $\abs{\operatorname{Re}(z)}\leq\abs{z}$ holds by preliminary square comparison and nonnegativity.
\item Set $q\defeq z\overline{w}$. The identity $\overline{q}=\overline{z}w$ holds by Proposition \ref{prop:complex-conjugation-identities}(ii) and the definition of conjugation.\\
The bound $\operatorname{Re}(q)\leq\abs{\operatorname{Re}(q)}\leq\abs{q}$ holds by the real estimate and (iv). The identity $\abs{q}=\abs{z}\abs{w}$ holds by (iii).\\
By Proposition \ref{prop:complex-conjugation-identities}(i), $\abs{z+w}^2=(z+w)(\overline{z}+\overline{w})$. By distributivity, $(z+w)(\overline{z}+\overline{w})=\abs{z}^2+q+\overline{q}+\abs{w}^2$.\\
Identity $q+\overline{q}=2\operatorname{Re}(q)$ holds by Proposition \ref{prop:complex-conjugation-identities}(iii). $\abs{z+w}^2\leq\abs{z}^2+2\abs{z}\abs{w}+\abs{w}^2=(\abs{z}+\abs{w})^2$ by preceding bounds.\\
$\abs{z+w}\leq\abs{z}+\abs{w}$ by the preliminary square comparison and nonnegativity.
\end{enumerate}
\hlt{Structural Commentary}: Squared moduli prove (ii)--(iv); the expansion in (v) reduces its cross term to (iii)--(iv).\\
\hlt{Key Insight}: The identity $z\overline{z}=a^2+b^2$ converts complex algebra into ordered real estimates.
\end{proof}
\end{theorem}
